\section*{Abstract}
\addcontentsline{toc}{section}{Abstract}

Digital circuits form the basis of computers and most modern electronic devices.
Building a physical prototype for each design change is slow and costly, so engineers test designs in software simulators before manufacture.
This project implements the circuit modeling part of the object-oriented approach proposed by Dezan, Lagadec, and Pottier as a simulator written in Java.
In this approach, software objects represent hardware elements such as logic gates and memory elements.

A user describes a circuit in a text file called a netlist.
The simulator calculates how the signals in the circuit change over time.
It presents the results as graphs called waveforms, as schematic diagrams, and as truth tables.
The system extends the original approach with signal delays, simulated hardware faults, timing analysis, and a measure of how many faults a test can detect.
It runs as a desktop application, a command line tool, and a browser application, and all three share one simulation engine.

An automated suite of 122 tests covers 91.3 percent of the source code lines outside the desktop interface.
This figure exceeds the 80 percent target that the project plan set for the domain logic and simulation engine.

\vspace{6pt}
\textbf{Keywords.} Digital circuit simulation, discrete-event simulation, object-oriented design, design patterns, fault injection.
